\documentclass[11pt]{article}
\usepackage[a4paper,margin=27mm]{geometry}
\usepackage[T1]{fontenc}
\usepackage{lmodern,microtype}
\usepackage{amsmath,amssymb,amsthm,mathtools}
\usepackage{enumitem,xcolor}
\usepackage[colorlinks=true,linkcolor=blue!55!black,citecolor=blue!55!black,urlcolor=blue!55!black]{hyperref}
\hypersetup{
 pdftitle={Exact Floors and Proper-Span Extremizers for Higher Osculating Absorption},
 pdfauthor={Lluis Eriksson},
 pdfsubject={Algebraic geometry; osculating spaces; jets; interpolation},
 pdfkeywords={osculating spaces, fat points, higher contact, tangent absorption, arbitrary characteristic}
}
\newtheorem{theorem}{Theorem}[section]
\newtheorem{proposition}[theorem]{Proposition}
\newtheorem{lemma}[theorem]{Lemma}
\newtheorem{corollary}[theorem]{Corollary}
\theoremstyle{definition}
\newtheorem{definition}[theorem]{Definition}
\theoremstyle{remark}
\newtheorem{remark}[theorem]{Remark}
\theoremstyle{plain}
\newcommand{\PP}{\mathbf P}
\newcommand{\kk}{\mathbf k}
\newcommand{\cO}{\mathcal O}
\newcommand{\cI}{\mathcal I}
\newcommand{\Span}{\operatorname{Span}}
\newcommand{\Osc}{\operatorname{Osc}}

\title{Exact Floors and Proper-Span Extremizers\\for Higher Osculating Absorption}
\author{Lluis Eriksson\\\small Independent researcher}
\date{August 24, 2026\quad (version 2)}

\begin{document}
\maketitle

\begin{abstract}
Let $X$ be a smooth projective integral $d$-fold over an algebraically
closed field, let $H$ be very ample, and use the complete embedding defined
by $H^m$.  Fix $1\leq s\leq m$.  A nonempty finite reduced set
$Z\subset X$ is point-span $s$-osculating-absorbing if the span $S_Z$ of
its value evaluations contains the affine osculating space of order $s$ at
every support.  We prove the exact characteristic-free floor
\[
 \dim S_Z,\ |Z|\geq \binom{d+m}{d}.
\]
The estimate follows from the exact tangent-absorption theorem, but its
sharpness under the stronger higher-order hypothesis is not formal.  For
every $(d,m,s)$ and every algebraically closed field we construct a smooth
integral hypersurface and a proper-span equality set with
\(\dim S_Z=|Z|=\binom{d+m}{d}\).  All supports lie in one hyperplane whose
normal coordinate vanishes to order $s+1$ on the hypersurface.  The proof
uses two explicit incidence estimates and is valid in positive
characteristic.

We also prove a rank-sensitive refinement.  If $r_1(Z)$ is the evaluation
rank in $H$ and $m\geq2s+1$, then
\[
 \dim S_Z,\ |Z|\geq
 \max\left\{\binom{d+m}{d},\binom{d+s}{d}r_1(Z)\right\}.
\]
The degree threshold is necessary for a uniform statement, as rational
normal curves show.  The earlier mixed-jet certificate is retained for
jet-ample polarizations, but its numerical floor is explicitly identified
as nonoptimal for tensor powers after the exact tangent theorem.  No equality
classification, minimal hypersurface degree, or exhaustive priority claim is
made.
\end{abstract}

\section{Statement and scope}

Let $\kk$ be algebraically closed, let $X$ be smooth, projective, integral,
and of dimension $d\geq1$, and let $H$ be very ample.  Write
\[
 \varphi_m:X\longrightarrow\PP(H^0(X,H^m)^*)
\]
for the complete embedding.  For a nonempty finite reduced set $Z\subset X$,
set
\[
 S_Z=\operatorname{Im}\!\left(
 H^0(Z,H^m|_Z)^*\longrightarrow H^0(X,H^m)^*\right).
\]
Equivalently, $S_Z$ is the span of the evaluation lines at the supports;
this formulation makes no choice of fiber trivializations.
For $p\in X$ and $a\geq0$, put
\[
 (a+1)p=\operatorname{Spec}(\cO_{X,p}/\mathfrak m_p^{a+1}),\qquad
 \lambda_d(a)=\binom{d+a}{d},
\]
and set $\lambda_d(-1)=0$.

The order-$a$ principal-parts evaluation is
\[
 j^a_p:H^0(X,H^m)\longrightarrow
 H^0\!\left((a+1)p,H^m|_{(a+1)p}\right).
\]
We define the affine order-$a$ osculating space intrinsically by
\[
 \widehat\Osc^a_p(H^m)=\operatorname{Im}((j^a_p)^*)
 \subset H^0(X,H^m)^*.
\]
This convention uses no ordinary derivatives or factorials.

\medskip

\begin{definition}
Fix $1\leq s\leq m$.  The set $Z$ is
\emph{point-span $s$-osculating-absorbing} for $H^m$ if
\[
 \widehat\Osc^s_p(H^m)\subseteq S_Z\qquad(p\in Z).
\]
\end{definition}

Put
\[
 N_{d,m}:=\binom{d+m}{d}.
\]

\begin{theorem}[Exact higher-osculating floor]\label{thm:main}
Over an algebraically closed field of arbitrary characteristic, every
nonempty finite reduced point-span $s$-osculating-absorbing set for the
complete $H^m$-embedding satisfies
\[
 \dim S_Z\geq N_{d,m},\qquad |Z|\geq N_{d,m}.
\]
\end{theorem}

\begin{theorem}[Rank-sensitive higher blocks]\label{thm:rank-sensitive}
Let
\[
 r_1(Z)=\operatorname{rk}\bigl(H^0(X,H)\longrightarrow
 H^0(Z,H|_Z)\bigr).
\]
If $m\geq2s+1$, then
\[
 \dim S_Z,\ |Z|\geq
 \max\left\{N_{d,m},\lambda_d(s)r_1(Z)\right\}.
\]
The threshold $m\geq2s+1$ cannot be weakened in a uniform theorem with the
same second term.
\end{theorem}

\begin{theorem}[Characteristic-free proper-span equality]
\label{thm:sharpness}
For every algebraically closed field, every $d,m\geq1$, and every
$1\leq s\leq m$, there exist a smooth integral hypersurface
$X^d\subset\PP^{d+1}$, a nonempty reduced set $Z\subset X$, and
$H=\cO_X(1)$ such that
\[
 |Z|=\dim S_Z=N_{d,m}
 <h^0(X,H^m)=\binom{d+m+1}{d+1},
\]
and $Z$ is point-span $s$-osculating-absorbing for $H^m$.  All supports may
lie in one hyperplane whose defining coordinate belongs to
$\mathfrak m_{X,p}^{s+1}$ for every $p\in Z$.
\end{theorem}

Theorem~\ref{thm:main} supersedes the smaller mixed floor published in the
first version of this record.  That earlier inequality remains correct, as
does its heterogeneous interpolation certificate, but higher-order
absorption contains tangent absorption and the latter now has the exact
binomial floor \cite{ErikssonExactFloor2026}.  The contributions of this
revision are
the higher-block refinement and proper-span sharpness in every order and
characteristic.

\section{The annihilator criterion}

Let $(s+1)Z=\coprod_{p\in Z}(s+1)p$.  Finite-scheme duality gives
\[
 S_Z^\perp=H^0(X,\mathcal I_Z\otimes H^m)
\]
and
\[
 \left(\sum_{p\in Z}\widehat\Osc^s_p(H^m)\right)^\perp
 =H^0(X,\mathcal I_{(s+1)Z}\otimes H^m).
\]

\begin{lemma}[Higher-order kernel criterion]\label{lem:kernel}
The set $Z$ is point-span $s$-osculating-absorbing if and only if
\[
 H^0(X,\mathcal I_Z\otimes H^m)
 =H^0(X,\mathcal I_{(s+1)Z}\otimes H^m).
\]
\end{lemma}

\begin{proof}
The inclusion from right to left always holds because
$\mathcal I_{(s+1)Z}\subseteq\mathcal I_Z$.  Absorption is equivalent to
\[
 \sum_{p\in Z}\widehat\Osc^s_p(H^m)\subseteq S_Z.
\]
Taking annihilators reverses this inclusion and supplies the converse
inclusion between the displayed spaces of sections.
\end{proof}

The principal-parts filtration gives
$\widehat\Osc^a_p(H^m)\subseteq\widehat\Osc^s_p(H^m)$ for $a\leq s$.
Thus order-$s$ absorption contains every lower-order jet block.  The converse
is not asserted.

\begin{proof}[Proof of Theorem~\ref{thm:main}]
Because $s\geq1$, every affine tangent space
$\widehat\Osc^1_p(H^m)$ is contained in
$\widehat\Osc^s_p(H^m)$ and hence in $S_Z$.  Thus $Z$ is point-span
tangent-absorbing.  The exact tangent-absorption theorem
\cite[Theorem~1.1]{ErikssonExactFloor2026} gives
\[
 \dim S_Z,\ |Z|\geq\binom{d+m}{d}=N_{d,m}.
\]
No characteristic-zero derivative or factorial is introduced by this
reduction.
\end{proof}

\begin{lemma}[Downward propagation]\label{lem:downward}
If $Z$ is point-span $s$-osculating-absorbing for $H^m$, then it is
point-span $s$-osculating-absorbing for $H^k$ for every $s\leq k\leq m$.
\end{lemma}

\begin{proof}
Fix $a\in H^0(X,\cI_Z\otimes H^k)$ and $p\in Z$.  Choose
$b_p\in H^0(X,H^{m-k})$ with $b_p(p)\ne0$; take $b_p=1$ if $k=m$.
Then $a b_p$ vanishes on $Z$, so Lemma~\ref{lem:kernel} in degree $m$
makes it vanish on $(s+1)p$.  The germ of $b_p$ is a unit, hence the germ
of $a$ belongs to $\mathfrak m_p^{s+1}$.  This holds at every support, so
the two kernels in degree $k$ agree and Lemma~\ref{lem:kernel} applies.
\end{proof}

\section{Mixed-jet interpolation}

\begin{lemma}[Full jets from powers]\label{lem:single}
If $p\in X$ and $0\leq a\leq k$, then
\[
 H^0(X,H^k)\longrightarrow
 H^0\!\left((a+1)p,H^k|_{(a+1)p}\right)
\]
is surjective.  Hence
$\dim\widehat\Osc^a_p(H^k)=\lambda_d(a)$.
\end{lemma}

\begin{proof}
Choose $u_0\in H^0(X,H)$ nonzero at $p$.  Very ampleness and smoothness give
$u_1,\ldots,u_d$, vanishing at $p$, such that $x_i=u_i/u_0$ form a regular
system of parameters.  For every multi-index $\alpha$ with $|\alpha|\leq a$,
\[
 u_0^{k-|\alpha|}u_1^{\alpha_1}\cdots u_d^{\alpha_d}
\]
has local representative $x^\alpha$ after trivializing by $u_0^k$.
These monomials form a basis modulo $\mathfrak m_p^{a+1}$.  No
differentiation is used.
\end{proof}

\begin{lemma}[Higher-block extension]\label{lem:higher-extension}
Assume $m\geq2s+1$.  Suppose the schemes
$(s+1)p_1,\ldots,(s+1)p_t$ impose independent conditions on $H^m$.
If the $H$-evaluation at $x\in X$ is not in the span of the evaluations at
$p_1,\ldots,p_t$, then adjoining $(s+1)x$ preserves independence.
\end{lemma}

\begin{proof}
Choose $e\in H^0(X,H)$ vanishing at every $p_i$ and nonzero at $x$.
Then $e^{s+1}$ vanishes on all old $(s+1)$-neighbourhoods and is a unit on
$(s+1)x$.  Put $k=m-s-1$.  The numerical hypothesis gives $k\geq s$, so
Lemma~\ref{lem:single} says that $H^k$ realizes arbitrary order-$s$ data at
$x$.  Multiplication by the local unit $e^{s+1}$ preserves that jet target
and kills all old targets.  Correcting residual data at $x$ proves
surjectivity on the enlarged disjoint union.
\end{proof}

\begin{proof}[Proof of Theorem~\ref{thm:rank-sensitive}]
Choose $r_1(Z)$ supports whose degree-one evaluation lines form a basis,
ordered so that each new line escapes the preceding span.  A single
$(s+1)$-neighbourhood imposes $\lambda_d(s)$ independent conditions on
$H^m$ by Lemma~\ref{lem:single}.  Repeated use of
Lemma~\ref{lem:higher-extension} makes the $r_1(Z)$ selected
neighbourhoods independent.  Their dual order-$s$ osculating blocks form a
direct sum of dimension $\lambda_d(s)r_1(Z)$ inside $S_Z$.  Combine this
with Theorem~\ref{thm:main} and $|Z|\geq\dim S_Z$.

To see that the degree threshold is necessary for this uniform second term,
take $X=\PP^1$, $H=\cO_{\PP^1}(1)$, and $m+1$ distinct points.  Their
$H^m$-evaluation span is the whole $(m+1)$-dimensional ambient vector space,
so they absorb every order $s\leq m$, while $r_1(Z)=2$.  If $m\leq2s$,
then $m+1<2(s+1)=\lambda_1(s)r_1(Z)$.
\end{proof}

\begin{lemma}[Mixed-jet interpolation]\label{lem:mixed}
Let $t\geq1$, let $p_1,\ldots,p_t$ be distinct points, let $a_i\geq0$, and put
\[
 K=\sum_{i=1}^t(a_i+1)-1.
\]
Then restriction is surjective:
\[
 H^0(X,H^K)\longrightarrow
 \bigoplus_{i=1}^t
 H^0\!\left((a_i+1)p_i,H^K|_{(a_i+1)p_i}\right).
\]
The same holds for $H^k$ in every degree $k\geq K$.
\end{lemma}

\begin{proof}
For each $i\ne j$, choose
$\ell_{ij}\in H^0(X,H)$ with
$\ell_{ij}(p_j)=0$ and $\ell_{ij}(p_i)\ne0$, and set
\[
 P_i=\prod_{j\ne i}\ell_{ij}^{a_j+1}.
\]
It is a unit at $p_i$ and vanishes to the prescribed order at every other
support.  Lemma~\ref{lem:single} realizes arbitrary order-$a_i$ data using
$H^{a_i}$.  Multiplication by the unit $P_i$ is an automorphism of the
truncated local algebra, so degree-$K$ sections realize arbitrary data at
$p_i$ and zero data elsewhere.  Sum over $i$.

For $k>K$, multiply by a section of $H^{k-K}$ nonzero at every support.
It exists because $\kk$ is infinite and a finite union of proper linear
subspaces cannot exhaust the section space.
\end{proof}

Over the complex numbers this is also the line-bundle specialization of
\cite[Proposition~2.3]{BDRS1999}.  The direct proof records the exact
specialization and is characteristic-free.  The interpolation statement itself
is not claimed as new; the point here is its rank consequence under the
absorption constraint and the ensuing optimization.

\section{The retained mixed-jet certificate}

Write uniquely
\[
 m+1=q(s+1)+r,\qquad q\geq1,\quad 0\leq r\leq s,
\]
and put
\begin{equation}\label{eq:B}
 B_{d,s}(m)=q\lambda_d(s)+\lambda_d(r-1).
\end{equation}

\begin{theorem}[Mixed-jet floor from version one]\label{thm:mixed-floor}
Under the hypotheses of Theorem~\ref{thm:main},
\[
 \dim S_Z,\ |Z|\geq B_{d,s}(m).
\]
For $d\geq2$ and $s<m$, this valid inequality is strictly weaker than the
exact floor $N_{d,m}$ of Theorem~\ref{thm:main}; it is retained because the
heterogeneous certificate below also applies to abstract jet-ample
polarizations.
\end{theorem}

The interpolation lemma first gives a heterogeneous certificate, before
optimization.

\begin{proposition}[Mixed-block rank certificate]\label{prop:certificate}
Under the hypotheses of Theorem~\ref{thm:main}, choose $t\geq1$ distinct supports
$p_1,\ldots,p_t\in Z$ and orders $0\leq a_i\leq s$.  If
\[
 \sum_{i=1}^t(a_i+1)\leq m+1,
\]
then
\[
 \dim S_Z\geq\sum_{i=1}^t\lambda_d(a_i).
\]
\end{proposition}

\begin{proof}
Put $K=\sum_i(a_i+1)-1\leq m$.  Lemma~\ref{lem:mixed}, followed by degree
raising when $K<m$, says that the selected infinitesimal neighborhoods impose
independent conditions on $H^m$.  Dually, their osculating blocks form a direct
sum of dimension $\sum_i\lambda_d(a_i)$.  Each block lies in $S_Z$ because
$a_i\leq s$ and $Z$ is order-$s$ absorbing.
\end{proof}

\begin{lemma}[Discrete convex packing]\label{lem:packing}
If $1\leq w_i\leq s+1$ and $\sum_iw_i\leq m+1$, then
\[
 \sum_i\binom{d+w_i-1}{d}\leq B_{d,s}(m).
\]
Equality is attained by $q$ weights $s+1$ and, if $r>0$, one weight $r$.
\end{lemma}

\begin{proof}
Put $f(0)=0$ and $f(w)=\binom{d+w-1}{d}$ for $w\geq1$.  Its forward differences
\[
 f(w+1)-f(w)=\binom{d+w-1}{d-1}
\]
are nondecreasing for $w\geq1$, while $f(1)-f(0)=1$.  Moving one unit from a
smaller block to a larger block below the cap $s+1$ does not decrease the
total; a weight-one block disappears when it loses its unit.  Repetition packs
the budget into full blocks and at most one residual block.  If the original total
weight is smaller than $m+1$, first enlarge a nonfull block or append a
weight-one block; either operation strictly increases the objective.  Thus an
optimum uses the whole budget, and the packed pattern gives \eqref{eq:B}.
\end{proof}

\begin{proof}[Proof of Theorem~\ref{thm:mixed-floor}]
Put $n=|Z|$.  If $r=0$ and $q=1$, nonemptiness gives $n\geq1=q$.  Now
suppose $r=0$, $q\geq2$, and $n\leq q-1$.  Then $n\geq1$, and the union of the full
order-$s$ neighborhoods at all supports is independent in degree
\[
 (s+1)n-1\leq(s+1)(q-1)-1<m
\]
by Lemma~\ref{lem:mixed}, and remains so after raising degree to $m$.
Its dual has dimension $n\lambda_d(s)$ and lies in $S_Z$ by absorption.
This contradicts $\dim S_Z\leq n$, since $\lambda_d(s)>1$.  Hence $n\geq q$.

If $r>0$ and $n\leq q$, the same contradiction applies because
\[
 (s+1)n-1\leq q(s+1)-1=m-r<m.
\]
Thus $n\geq q+1$.

Choose $q$ supports with full order-$s$ blocks and, if $r>0$, one further
support with an order-$(r-1)$ block.  Their mixed interpolation degree is
\[
 q(s+1)+r-1=m.
\]
Their independent dual spaces lie in $S_Z$ and have total dimension
$B_{d,s}(m)$.  Therefore $\dim S_Z\geq B_{d,s}(m)$, and
$|Z|\geq\dim S_Z$ because $S_Z$ is generated by $|Z|$ value lines.
\end{proof}

The certificate and packing argument do not intrinsically require a tensor
power.  The following form records the more general input precisely.

\begin{corollary}[Jet-ample polarization]\label{cor:jetample}
Let $L$ be an $M$-jet ample line bundle on $X$, meaning that for every set of
distinct points $p_i$ and positive integers $b_i$ with
$\sum_i b_i=M+1$, the restriction map
\[
 H^0(X,L)\longrightarrow
 \bigoplus_i H^0\!\left(b_i p_i,L|_{b_i p_i}\right)
\]
is surjective.  Fix $1\leq s\leq M$, write
$M+1=Q(s+1)+R$ with $0\leq R\leq s$, and define $S_Z$ and the osculating
spaces using $L$ in place of $H^m$.  If a nonempty finite reduced set $Z$ is
point-span $s$-osculating-absorbing, then
\[
 \dim S_Z,\ |Z|\geq
 Q\binom{d+s}{d}+\binom{d+R-1}{d},
\]
where the last summand is interpreted as zero when $R=0$.
Over the complex numbers, repeated application of
\cite[Proposition~2.3]{BDRS1999} gives this conclusion with $M=\ell a$ for
$L=A^{\ell}$ whenever $A$ is $a$-jet ample and $1\leq s\leq\ell a$.
Here $a\geq1$ and $\ell\geq1$ are integers.
\end{corollary}

\begin{proof}
The stated jet-ampleness also gives surjectivity when $\sum_i b_i<M+1$:
add one point outside the selected finite set with the missing positive
weight, and then project away its jet target.  Such a point exists because
$d\geq1$ and the algebraically closed field is infinite.  Thus every selected
collection with $b_i=a_i+1\leq s+1$ and total weight at most $M+1$ is
independent.  The support-threshold proof above applies verbatim with $M$ in
place of $m$, including the separate case $Q=1,R=0$.  Selecting $Q$ full
blocks and, when $R>0$, one residual block gives the displayed rank.  The
last assertion is the tensor-product theorem cited above, iterated $\ell-1$
times.
\end{proof}

\section{Characteristic-free unisolvent supports}

\begin{lemma}[Existence of an evaluation basis]\label{lem:unisolvent}
For every algebraically closed field and every $d,m\geq1$, there is a set of
$N_{d,m}$ distinct points $Z\subset\PP^d$ such that evaluation is an
isomorphism
\[
 H^0(\PP^d,\cO(m))\xrightarrow{\ \sim\ }H^0(Z,\cO_Z(m)).
\]
\end{lemma}

\begin{proof}
The evaluation lines of all $\kk$-points span
$H^0(\PP^d,\cO(m))^*$.  Otherwise their annihilator would contain a nonzero
homogeneous degree-$m$ form vanishing at every $\kk$-point of $\PP^d$,
which is impossible over the infinite field $\kk$.  Select a basis from
these evaluation lines.  Very ampleness separates points, so their
$N_{d,m}$ supports are distinct; dualizing gives the displayed isomorphism.
\end{proof}

This existence proof replaces the integer simplex lattice used in version
one.  That lattice remains a useful characteristic-zero fixture, but its
points can collide and its falling-factorial basis can degenerate when the
characteristic is at most $m$.

\section{Smooth higher-contact extremizers}

Let $W=V(y)\cong\PP^d$ be a hyperplane in $\PP^{d+1}$ and let
$Z\subset W$ be an unisolvent set from Lemma~\ref{lem:unisolvent}.  Write
$(s+1)Z_W$ for the disjoint union of the subschemes defined by
$\mathfrak m_{W,p}^{s+1}$; equivalently,
\[
 \cI_{(s+1)Z_W}=\bigcap_{p\in Z}\mathfrak m_{W,p}^{s+1}.
\]
This is a fat-point or symbolic-power ideal sheaf, not an assertion about
the ordinary power of the homogeneous ideal of $Z$.

\begin{proposition}[Smooth prescribed higher contact]
\label{prop:higher-contact}
Put $N=N_{d,m}$ and $E=(s+1)N$.  For every $n\geq E+1$, there is a smooth
integral degree-$n$ hypersurface $X\subset\PP^{d+1}$ containing $Z$ such
that
\[
 y\in\mathfrak m_{X,p}^{s+1}\qquad(p\in Z).
\]
In particular $W$ is the tangent hyperplane at every support and has contact
through order $s$ there.
\end{proposition}

\begin{proof}
Set
\[
 V=H^0\bigl(W,\cI_{(s+1)Z_W}(n)\bigr).
\]
We first choose $f\in V$ whose zero divisor is smooth on $W\setminus Z$.
Fix $q\in W\setminus Z$.  For every $p\in Z$, choose a hyperplane
$\ell_{p,q}$ through $p$ and not through $q$, and put
\[
 A_q=\prod_{p\in Z}\ell_{p,q}^{s+1}.
\]
Then $A_q$ has degree $E$, belongs to
$H^0(W,\cI_{(s+1)Z_W}(E))$, and is a unit on the first neighbourhood $2q$.
Since $n-E\geq1$, $\cO_W(n-E)$ generates all first jets at $q$.
Multiplication by $A_q|_{2q}$ therefore proves that
\[
 V\longrightarrow H^0(2q,\cO_{2q}(n))
\]
is surjective.

Consider the incidence
\[
 \Sigma_1=\{(q,[f])\in(W\setminus Z)\times\PP(V):j_q^1f=0\}.
\]
For fixed $q$, vanishing of the first jet has codimension $d+1$; hence
\[
 \dim\Sigma_1\leq d+\dim\PP(V)-(d+1)
 =\dim\PP(V)-1.
\]
The closure of its image cannot fill $\PP(V)$.  Choose $f$ outside that
closure.  Then $V(f)$ is smooth at every one of its points outside $Z$.
This is an incidence calculation with local jets, so no Euler identity,
ordinary derivative convention, or characteristic restriction is involved.

Fix such an $f$ and let
\[
 U=H^0(\PP^{d+1},\cO(n-1)),\qquad F_G=f+yG\quad(G\in U).
\]
On the open set $y\ne0$, multiplication by $y$ is an isomorphism on first
neighbourhoods.  Since $n-1\geq1$, the affine map
\[
 U\longrightarrow H^0(2q,\cO_{2q}(n)),\qquad
 G\longmapsto j_q^1(f+yG)
\]
is surjective for every such $q$.  The condition that $F_G$ be singular at
a fixed point of the $(d+1)$-dimensional ambient open set has codimension
$d+2$.  Thus
\[
 \Sigma_2=\{(q,G):y(q)\ne0,\ j_q^1F_G=0\}
\]
has dimension at most
$(d+1)+\dim U-(d+2)=\dim U-1$.  The closure of its image is therefore a
proper closed subset of $U$.
The additional conditions $G(p)\ne0$, $p\in Z$, are the complements of
finitely many hyperplanes in the irreducible affine space $U$.  Choose $G$
outside the closure of the incidence image and outside all those hyperplanes.

Let $X=V(F_G)$.  It is smooth on $y\ne0$ by the second incidence.  If
$q\in W\setminus Z$ lies on $X$, then $f(q)=0$ and the tangential first jet
of $f$ is nonzero, so $X$ is smooth at $q$.  At $p\in Z$, all first jets of
$f$ vanish and the normal first jet of $F_G$ is $G(p)\ne0$, so $X$ is
smooth there as well.

The hypersurface is connected.  Indeed, from
\[
 0\longrightarrow\cO_{\PP^{d+1}}(-n)
 \longrightarrow\cO_{\PP^{d+1}}
 \longrightarrow\cO_X\longrightarrow0
\]
and $H^1(\PP^{d+1},\cO(-n))=0$ for $d+1\geq2$, one gets
$H^0(X,\cO_X)=\kk$.  A smooth scheme has disjoint irreducible components;
connectedness therefore makes $X$ integral.

Finally, regard $f$ as its lift independent of $y$.  Since that lift belongs
to $\mathfrak m_{\PP^{d+1},p}^{s+1}$, its image belongs to
$\mathfrak m_{X,p}^{s+1}$.  In $\cO_{X,p}$ the germ of $G$ is a unit and
\[
 y=-f/G\in\mathfrak m_{X,p}^{s+1},
\]
which proves the required contact.
\end{proof}

\begin{proof}[Proof of Theorem~\ref{thm:sharpness}]
Choose $Z\subset W$ by Lemma~\ref{lem:unisolvent}, take
$n\geq(s+1)N_{d,m}+1$, and choose $X$ by
Proposition~\ref{prop:higher-contact}.  In particular $n>m$.  The
hypersurface sequence twisted by $m$, together with
\[
 H^0(\PP^{d+1},\cO(m-n))=0,
 \qquad H^1(\PP^{d+1},\cO(m-n))=0,
\]
gives an isomorphism
\begin{equation}\label{eq:ambient-restriction}
 H^0(\PP^{d+1},\cO(m))\xrightarrow{\ \sim\ }
 H^0(X,\cO_X(m)).
\end{equation}

Let a section on the right vanish on $Z$ and denote its unique ambient lift
by $Q$.  Its restriction $Q|_W$ is a degree-$m$ form vanishing at the
unisolvent set $Z$, hence $Q|_W=0$.  Thus $Q=yR$ for a degree-$(m-1)$ form
$R$.  Proposition~\ref{prop:higher-contact} gives
\[
 Q=yR\in\mathfrak m_{X,p}^{s+1}\qquad(p\in Z).
\]
Therefore
\[
 H^0(X,\cI_Z\otimes\cO_X(m))
 =H^0(X,\cI_{(s+1)Z}\otimes\cO_X(m)),
\]
and Lemma~\ref{lem:kernel} proves order-$s$ absorption.

Unisolvence gives $\dim S_Z=|Z|=N_{d,m}$.  Equation
\eqref{eq:ambient-restriction} gives
\[
 h^0(X,\cO_X(m))=\binom{d+m+1}{d+1}>N_{d,m},
\]
so the point span is proper.
\end{proof}

\begin{corollary}[Exact universal minimum]\label{cor:minimum}
For every algebraically closed field and every $d,m\geq1$,
$1\leq s\leq m$, the minimum possible cardinality and span dimension of a
point-span $s$-osculating-absorbing set, as $(X,H,Z)$ vary under the standing
hypotheses, are both $N_{d,m}$.  The minimum remains the same after requiring
the point span to be proper.
\end{corollary}

\begin{proof}
The lower bound is Theorem~\ref{thm:main}; the proper-span equality examples
are Theorem~\ref{thm:sharpness}.
\end{proof}

The degree $n\geq(s+1)N_{d,m}+1$ is a transparent sufficient threshold,
not a minimality claim.  The proof certifies an open set of choices rather
than one distinguished equation.

\section{Exact computational witnesses}

The accompanying replay uses exact rational and finite-field linear algebra.
It checks the original mixed-block identities, the new rank-sensitive formula
and its $\PP^1$ falsification boundary, unisolvent evaluation sets in several
small characteristics, and local fixtures in which the normal coordinate lies
in $\mathfrak m^{s+1}$.  Fresh JSON output is compared byte for byte with the
committed result.  These are finite diagnostic witnesses: they do not prove
the incidence dimension counts, smoothness of a general member, the universal
rank theorem, or priority.

\section{Context and limitations}

Jet ampleness encodes simultaneous interpolation on fat points.  A
tensor-product generation result over the complex numbers appears in
\cite{BDRS1999}.  Higher osculating spaces and fundamental forms are
formulated through principal-parts sheaves over algebraically closed fields
in \cite{MallavibarrenaPiene2024}; that is the intrinsic language used here.

Ballico's $X$-rank of a linear subspace is the direct established language
for the minimum number of points of $X$ whose span contains a prescribed
linear space \cite{BallicoXRank2013}.  That work computes ranks of
tangent-containing subspaces for rational normal curves and several
Veronese configurations.  The present condition is simultaneous and
self-referential: one reduced set must absorb the order-$s$ osculating space
at every one of its own supports.  We do not claim to introduce subspace rank
or point-span absorption as a general concept.

Osculating spaces, higher Gauss maps, and secant varieties have a substantial
literature \cite{BallicoEtAl2004,BallicoFontanari2004,BernardiEtAl2009}.
The Veronese literature relates osculating secant defectivity to Hilbert
functions of fat points.  Those questions differ from proper-span equality
for simultaneous absorption, but they delimit the novelty claim.  This
selective comparison is not a priority certification.

The authorial note B215 contains the order-two mixed interpolation pattern
\cite{ErikssonB215}; it is an antecedent, not independent validation.
Version one of this ARR record extracted that certificate and optimized its
packing formula.  The later exact tangent theorem
\cite{ErikssonExactFloor2026} made that numerical floor nonoptimal for tensor
powers.  This major revision records the corrected hierarchy rather than
concealing the chronology.

The limitations are material.
\begin{itemize}
\item The higher-order absorption hypothesis is explicit.  Tangent absorption
alone is not claimed to imply second- or higher-order absorption.
\item The restriction $s\leq m$ ensures the full local rank in
Lemma~\ref{lem:single}.  No formula for $s>m$ is asserted.
\item Completeness of the $H^m$ system is used because all separator products
must be available.  Arbitrary subsystems are outside the theorem.
\item The higher-block term requires $m\geq2s+1$.  The curve counterexample
shows this threshold is necessary for that uniform formula, but no optimal
replacement is claimed below the threshold.
\item The extremizers prove existence at the sufficient hypersurface degree
$n\geq(s+1)N_{d,m}+1$.  They neither minimize $n$, classify equality sets,
nor provide one canonical equation.
\item The exact floor is invoked from a separate, earlier ARR record.  The
proper-span construction and higher-block proof are given here in full.
\item The literature search is selective.  The replay is not a proof checker,
and no independent human peer review or exhaustive priority review is
claimed.
\end{itemize}

\section{Conclusion}

Higher-osculating absorption has the same exact universal binomial floor as
tangent absorption, but the stronger condition is not geometrically empty.
Arbitrary-characteristic hypersurfaces with prescribed order-$(s+1)$ normal
contact realize the floor with a proper point span in every dimension,
degree, and allowed osculating order.  Away from that high-contact branch,
the rank-sensitive block argument forces
$\lambda_d(s)r_1(Z)$ independent conditions once $m\geq2s+1$.  The mixed-jet
certificate remains useful for abstract jet-ample polarizations, while its
older tensor-power floor is now explicitly subordinate to the exact theorem.
Minimal construction degree, equality classification, and stronger bounds
below the block threshold remain open.

{\small
\begin{thebibliography}{99}
\bibitem{BallicoEtAl2004}
E.~Ballico, C.~Bocci, E.~Carlini, and C.~Fontanari,
\emph{Osculating spaces to secant varieties},
arXiv:math/0406322 (2004).
\url{https://arxiv.org/abs/math/0406322}

\bibitem{BallicoFontanari2004}
E.~Ballico and C.~Fontanari,
\emph{On the osculatory behaviour of higher dimensional projective
varieties}, arXiv:math/0406319 (2004).
\url{https://arxiv.org/abs/math/0406319}

\bibitem{BallicoXRank2013}
E.~Ballico,
\emph{On the $X$-rank of linear subspaces},
Rev. Roumaine Math. Pures Appl. \textbf{58} (2013), 437--451.
\url{https://imar.ro/journals/Revue_Mathematique/pdfs/2013/4/4.pdf}

\bibitem{BDRS1999}
M.~C.~Beltrametti, S.~Di Rocco, and A.~J.~Sommese,
\emph{On generation of jets for vector bundles},
Rev. Mat. Complut. \textbf{12} (1999), 27--45.
\url{https://doi.org/10.5209/rev_REMA.1999.v12.n1.17182}

\bibitem{BernardiEtAl2009}
A.~Bernardi, M.~V.~Catalisano, A.~Gimigliano, and M.~Id\`a,
\emph{Secant varieties to osculating varieties of Veronese embeddings of
$\PP^n$}, J. Algebra \textbf{321} (2009), 982--1004.
\url{https://doi.org/10.1016/j.jalgebra.2008.10.020}

\bibitem{ErikssonExactFloor2026}
L.~Eriksson,
\emph{The exact rank floor for point-span tangent absorption in arbitrary
characteristic},
ARR-2026-2MHNZRRJP49Y9SWP, v3 (2026).
\url{https://arr-research.github.io/papers/ARR-2026-2MHNZRRJP49Y9SWP/versions/v3/}

\bibitem{ErikssonB215}
L.~Eriksson,
\emph{B215---Simultaneous second-jet interpolation raises the floor},
public working note (2026),
\href{https://github.com/lluiseriksson/hodge-conjecture-research-front/blob/main/proofs/B215-simultaneous-second-jet-interpolation-floor.md}{repository record B215}.

\bibitem{MallavibarrenaPiene2024}
R.~Mallavibarrena and R.~Piene,
\emph{On fundamental forms and osculating bundles},
Rend. Istit. Mat. Univ. Trieste \textbf{56} (2024), Art.~9, 21 pp.
\url{https://doi.org/10.13137/2464-8728/36877}
\end{thebibliography}
}

\end{document}
